% =====================================================================
%  Section 3 -- The differential systems and their local data
%  reader-integrity sentinel — non-ASCII, do not remove: é ö ±
%  CT-S37-DRAFT-B.  Draft only.
%
%  Every numerical value in this section carries the claim id it comes
%  from, and every epistemic label is the label that claim currently
%  carries in the corpus, checked by scripts/draft_check.py.  No figure
%  is written from memory or from a prior draft.
% =====================================================================

\section{The differential systems and their local data}\label{sec:systems}

\subsection{From the recurrence to System \texorpdfstring{$\mathsf A$}{A}}
\label{ssec:sysA}

Take $b(n)=3n^2+n+1$, the quadratic of the companion record, and let
$f(z)=\sum_{n\ge0}Q_nz^n$ be the generating function of the denominators
\eqref{eq:n-rec}. Then $f$ satisfies the inhomogeneous equation
\begin{equation}\label{eq:sysA}
  3z^3f''+10z^2f'+(z^2+5z-1)f=-1 ,
\end{equation}
which we call \emph{System $\mathsf A$}. This is \textbf{PROVEN}
(\texttt{S4-03}): the residual of \eqref{eq:sysA} against the exact integer
sequence $Q=1,5,76,2361,125209,10144290,\dots$ vanishes identically at every
order $z^0,\dots,z^{38}$ in exact arithmetic.

The homogeneous part of \eqref{eq:sysA} has projective normal potential
\begin{equation}\label{eq:rA}
  r_{\mathsf A}(z)=\frac{-3z^2-5z+3}{9z^3},
\end{equation}
and $z=0$ is its only finite singular point (\textbf{PROVEN}, \texttt{S4-05};
that no finite affine singular candidate exists outside $z=0$ is separately
\textbf{PROVEN} at \texttt{S6-C8}).

\subsection{Systems \texorpdfstring{$\mathsf B$}{B} and
\texorpdfstring{$\mathsf C$}{C}}\label{ssec:sysBC}

\emph{System $\mathsf B$} is the scalar reduction $L_\varphi$, whose
transcription is certified two independent ways (\textbf{PROVEN},
\texttt{S4-04}, and again at \texttt{S6-C3}): the factorisation
$q_2=4(2\sqrt3-3)z^2(z^2+z+3)$ expands exactly to the listed $q_2$, and
recomputing the normal potential $R=\tfrac14P^2+\tfrac12P'-Q$ of
\S\ref{ssec:objects}, here at $(P,Q)=(a,b)$ for the operator written
$\partial^2+a\partial+b$, from $(q_0,q_1,q_2)$ reproduces
\begin{equation}\label{eq:rB}
  r_{\mathsf B}(z)=\frac{11z^4+4z^2+4z+12}{4z^4(z^2+z+3)^2}
\end{equation}
exactly.

\emph{System $\mathsf C$} is the scalar equation
\begin{equation}\label{eq:sysC}
  (3x^2+x+1)\,W''+(6x+1)\,W'-x^2\,W=0 .
\end{equation}
It is in self-adjoint form: with $a(x)=3x^2+x+1$ one has $b(x)=6x+1=a'(x)$
exactly, so \eqref{eq:sysC} is $\bigl(a(x)W'\bigr)'+c(x)W=0$
(\textbf{PROVEN}, \texttt{S6-C12}; equivalently $b-a'=0$, \textbf{PROVEN},
\texttt{CT-S9-PVPARAMS-C005}). Its normal potential is
\begin{equation}\label{eq:rC}
  r_{\mathsf C}(x)=\frac{12x^4+4x^3+4x^2+11}{4\,(3x^2+x+1)^2}
\end{equation}
(\textbf{PROVEN}, \texttt{CT-S9-PVPARAMS-C009} and \texttt{S16-06}), and its
finite singular points are
\begin{equation}\label{eq:spm}
  s_\pm=\frac{-1\pm i\sqrt{11}}{6},
  \qquad \delta=s_--s_+=-\frac{i\sqrt{11}}{3}
\end{equation}
(\textbf{PROVEN}, \texttt{CT-S9-PVPARAMS-C006} and \texttt{S16-06}). Under
$x=s_++\delta z$ the potential \eqref{eq:rC} has poles only at $z=0$ and $z=1$,
both double with coefficient $-\tfrac14$, and no apparent singularity, so
$\mathsf C$'s three singular points are exactly $0,1,\infty$ after the
normalising M\"obius map (\textbf{PROVEN}, \texttt{S16-06}).

\subsection{Exact formal types}\label{ssec:formaltypes}

All three formal types below are exact, not asymptotic, and each is
\textbf{PROVEN}.

\paragraph{System $\mathsf A$ (\texttt{S4-05}).} At $z=0$ the pole order of
$r_{\mathsf A}$ is $3$, which is odd, so the singularity is irregular and
\emph{ramified}, with ramification index $2$, Katz invariant $\tfrac12$, and
exponential parts
\[
  \exp\Bigl(\mp\tfrac{2}{\sqrt3}\,z^{-1/2}\Bigr).
\]
At $z=\infty$ the pole order is again $3$: irregular, ramified, Katz invariant
$\tfrac12$.

\paragraph{System $\mathsf B$ (\texttt{S4-06}).} At $z=0$ the pole order is
$4$, which is even, so the singularity is irregular and \emph{unramified}, of
Poincar\'e rank $1$ and Katz invariant $1$, with exponential parts
$\exp(\mp(1/\sqrt3)/z)$. At $\rho_\pm=(-1\pm i\sqrt{11})/2$ the pole order is
$2$: regular singular. At $z=\infty$ the pole order is $2$: regular singular.

\paragraph{System $\mathsf C$ (\texttt{S6-C15}, \texttt{CT-S9-PVPARAMS-C011},
\texttt{CT-S9-PVPARAMS-C012}).} At infinity, in the coordinate $\tau=1/x$, the
normal-form potential has pole order $4$, leading coefficient $\tfrac13$, no
ramification, and Katz invariant --- equivalently Poincar\'e rank --- equal to
$1$. The scalar normal-form exponential factors there begin
\[
  \exp\Bigl(\tfrac{\sqrt3}{3}x\Bigr),\qquad
  \exp\Bigl(-\tfrac{\sqrt3}{3}x\Bigr),
\]
equivalently $\exp(-(\sqrt3/3)/\tau)$ and $\exp((\sqrt3/3)/\tau)$.

\paragraph{Actions read off the equations alone (\texttt{S4-07}).} No
convention has to be interpreted to extract the following, and it is
\textbf{PROVEN}. In the uniformiser $u=\sqrt z$, System $\mathsf A$ has single
action $2/\sqrt3=\xi_0$ and action \emph{difference} $4/\sqrt3$. System
$\mathsf B$, in $z$, has single action $1/\sqrt3$ and action difference
$2/\sqrt3=\xi_0$. The quantity $4/\sqrt3$ is therefore an action difference and
not the image of $\xi_0$ under a change of variables.

\subsection{The two finite singular points are not apparent}
\label{ssec:nonapparent}

This subsection is the one the companion record got wrong, and the correction
is stated here as mathematics; the history of the correction is not this
section's business.

At both finite singular points $s_\pm$ of System $\mathsf C$ the indicial
polynomial is $\lambda^2$, so the scalar local exponents are $\{0,0\}$ and the
exponent difference is $0$ (\textbf{PROVEN},
\texttt{CT-S9-PVPARAMS-C007}). Exponents $\{0,0\}$ do not by themselves decide
apparency: what decides it is the logarithm. Normalising $y_1(s)=1$, the second
Frobenius solution has \emph{nonzero} log coefficient
\begin{equation}\label{eq:logcoef}
  \frac{\sqrt{11}\,i}{11}\ \text{ at } s_-=\frac{-1-i\sqrt{11}}{6},
  \qquad
  -\frac{\sqrt{11}\,i}{11}\ \text{ at } s_+=\frac{-1+i\sqrt{11}}{6}
\end{equation}
(\textbf{PROVEN}, \texttt{CT-S9-PVPARAMS-C008}). Consequently \emph{neither
finite point is apparent}: at $s_-$ this is \textbf{PROVEN} at
\texttt{S6-C13} and at $s_+$ at \texttt{S6-C14}, in both cases with exact
repeated-root log residue $1$, and in both cases against the companion's
recorded exponents $\{0,1\}$.

In the normalised coordinate the same statement reads: the coefficient of
$(x-s)^{-2}$ in the normal-form potential is $-\tfrac14$ at both points, so the
normal-form exponent is repeated at $\tfrac12$ (\textbf{PROVEN},
\texttt{CT-S9-PVPARAMS-C010}).

The corresponding statements for System $\mathsf B$ are \textbf{PROVEN} in the
same way: $\rho_-$ at \texttt{S6-C4} and $\rho_+$ at \texttt{S6-C5}, both with
indicial equation $\lambda^2$, exponents $\{0,0\}$ and nonzero exact
repeated-root log residue; and $\infty$ at \texttt{S6-C6}, where the local
exponents are $1-\sqrt3/18$ and $2-\sqrt3/18$ and so are not integers.

Finally, the general assertion that such points are always apparent is refuted
on its own tested instances: for $b=x+1$ and $b=x^2+1$ the simple-root tests
give exponents $\{0,0\}$ with nonzero log residue $1$, so those instances are
\emph{not} apparent (\textbf{PROVEN}, \texttt{S6-C17}). That the companion's
$\beta=\gamma=0$ was read as reflecting apparency of both finite points is
\textbf{not supported} by the exact test --- a \textbf{STRUCTURAL} conclusion
(\texttt{S6-C19}), and stated with that label because it is a statement about
what an argument does and does not support.

\subsection{The Katz mismatch in a fixed coordinate, and the cover that
removes it}\label{ssec:cover}

Within a fixed base coordinate the formal types of $\mathsf A$ and $\mathsf B$
at the origin are incompatible: Katz $\tfrac12$ ramified against Katz $1$
unramified, and no gauge transformation and no unramified change of variable
can reconcile them. The ramified double cover
\begin{equation}\label{eq:cover}
  \pi:\ w\longmapsto z=w^2
\end{equation}
--- which is the uniformiser $u=\sqrt z$ of \S\ref{ssec:formaltypes} --- doubles
the Katz invariant and carries $\mathsf A$'s $\tfrac12$ to $\mathsf B$'s $1$.
The verdict drawn from this is \textbf{STRUCTURAL} and explicitly
\emph{inconclusive} as to global equivalence (\texttt{S4-09}).

The pullback is carried out through the normal-form rule
$R=(\varphi')^2\,r(\varphi)-\tfrac12\,\mathcal S(\varphi)$, with $\mathcal S$
the Schwarzian, which passes the Euler exponent-difference doubling test
(\textbf{STRUCTURAL}, \texttt{S5B-C2}); the CT-S4 baseline normal forms and
certificates recompute exactly beforehand (\textbf{PROVEN},
\texttt{S5B-C1}). Its \emph{vacuity pre-check} is \textbf{PROVEN}
(\texttt{S5B-C3}): under $z=w^2$ the model potential $r=1/z^4$ remains
unramified, while System $\mathsf A$'s pole order at the origin moves from $3$
to $4$ --- so the observable does move when the data move, and the cover is not
being credited with an effect it does not have.

After the pullback, $\pi^*\mathsf A$ has exactly one finite singular point
$w=0$, of pole order $4$, unramified, Katz invariant $1$ (\textbf{PROVEN},
\texttt{S5B-C4}), and an unramified irregular singularity of Katz invariant $1$
at infinity (\textbf{PROVEN}, \texttt{S5B-C5}).

Three statements bound what this establishes, and all three are
\textbf{STRUCTURAL}: the \emph{local} mismatch at the branch point is removed
(\texttt{S5B-C8}); the \emph{global} singular configurations of $\pi^*\mathsf A$
and $\mathsf B$ are \emph{not} reconciled by that cover alone
(\texttt{S5B-C9}); and CT-S4's verdict is therefore sharpened but not settled
as a global isomorphism (\texttt{S5B-C10}). The paper claims the local
statement and not the global one.

\subsection{Effective Katz multisets}\label{ssec:multisets}

Apparent singularities may be removed before a Katz multiset is read off, so
the multiset of record is the \emph{effective} one, computed after the
apparency tests of \S\ref{ssec:nonapparent} have been applied.

For System $\mathsf C$ the tests remove no finite points, and the effective
multiset is $(0,0,1,-)$, which matches van der Put--Saito's Table~2.1 row
$\widetilde D_5$/PV on p.~2635. This is \textbf{STRUCTURAL}
(\texttt{S6-C16}); the location of that row and of the selection rule on
p.~2634 is itself a reading of the printed page, \textbf{INSPECTED}
(\texttt{S6-C1}).

For System $\mathsf B$ the tests likewise remove no points, and the effective
multiset is $(1,0,0,0)$, which does \emph{not} match that row --- also
\textbf{STRUCTURAL} (\texttt{S6-C7}).

For System $\mathsf A$ the dash in $(\tfrac12,-,\tfrac12,-)$ denotes an
\emph{absent} singular point in the classification tuple, not a point dropped
as regular or as apparent (\textbf{PROVEN}, \texttt{S6-C9}).

\begin{remark}[A matching multiset is necessary, never sufficient]
\label{rem:necessary}
Table~2.1 classifies \emph{families} of connections. A single fixed operator
whose multiset matches a row has satisfied a necessary condition and nothing
more; family membership is not established until the deformation parameter is
specified. This is \textbf{STRUCTURAL} (\texttt{S6-C10}), and it is the reason
\S\ref{ssec:multisets} reports a match for $\mathsf C$ without assigning
$\mathsf C$ a Painlev\'e label. The label is not assigned anywhere in this
paper on the strength of a multiset.
\end{remark}

\subsection{\texorpdfstring{$\mathsf B$ and $\mathsf C$ under $z\mapsto1/x$}%
{B and C under z to 1/x}}\label{ssec:BC}

Under $z=1/x$, System $\mathsf B$ has \emph{exactly} the same normal-form
potential as System $\mathsf C$. The two original scalar equations are
nevertheless not the same operator: they differ by the explicit gauge factor
\begin{equation}\label{eq:gauge}
  x^{\sqrt3/18-1}\exp\Bigl(-\tfrac{\sqrt3}{3}x\Bigr),
\end{equation}
so they are gauge-equivalent in normal form without being literally equal.
This is \textbf{PROVEN} (\texttt{S6-C18}). Consistently, the inversion
$z\mapsto1/x$ applied to $\mathsf B$ makes $x=0$ ordinary, leaves finite
regular singular points at $(-1\pm i\sqrt{11})/6$, and gives an unramified Katz
invariant $1$ at infinity (\textbf{PROVEN}, \texttt{S5B-C6}).

That the same configuration was recorded independently in an earlier verdict
file is \textbf{INSPECTED} only, and with a stated gap: the apparency indices
were not recomputed there (\texttt{S5B-C7}). Nothing in this section rests on
that reading.

\medskip
\noindent
The distinction \eqref{eq:gauge} draws --- equal normal potential, different
operator --- is exactly the projective ambiguity that the later section on
removing that ambiguity resolves.
